% Original: PDF strana 18, gornji slajd.
\slajdid{02-s035}
\begin{frame}[label=02-s035]{Karnoove karte za tri promenljive}
\setlength{\parskip}{0pt}
% element: 02-s035:e01
\begin{columns}[T,onlytextwidth]
\begin{column}{.47\textwidth}\centering
\Oznaka{Dve kolone i četiri reda}\vspace{1mm}
\begingroup
\fontsize{14.52}{16.13}\selectfont
\renewcommand{\small}{\fontsize{14.52}{16.13}\selectfont}
\scalebox{.62}{%
\begin{karnaugh-map}(label=middle)[2][4][1][$A$][$B$][$C$]
\foreach \x/\y/\lab in {1/3/0,2/3/1,1/2/2,2/2/3,1/1/6,2/1/7,1/0/4,2/0/5}{
 \node[anchor=south east,inner sep=1.1pt,font=\fontsize{11.29}{12.9}\selectfont] at (\x,\y) {\lab};
}
\end{karnaugh-map}}
\endgroup

\end{column}
\begin{column}{.47\textwidth}\centering
\Oznaka{Četiri kolone i dva reda}\vspace{1mm}
\begingroup
\fontsize{14.52}{16.13}\selectfont
\renewcommand{\small}{\fontsize{14.52}{16.13}\selectfont}
\scalebox{.62}{%
\begin{karnaugh-map}(label=middle)[4][2][1][$A$][$B$][$C$]
\foreach \x/\y/\lab in {1/1/0,2/1/1,3/1/3,4/1/2,1/0/4,2/0/5,3/0/7,4/0/6}{
 \node[anchor=south east,inner sep=1.1pt,font=\fontsize{11.29}{12.9}\selectfont] at (\x,\y) {\lab};
}
\end{karnaugh-map}}
\endgroup

\end{column}
\end{columns}
\par
% element: 02-s035:e02
U levoj karti susedna su i polja \(0\leftrightarrow4\) i \(1\leftrightarrow5\). To možemo zamisliti kao presavijanje tabele oko horizontalne ose po sredini i spajanje gornje i donje ivice u prsten.
\par\vspace{1mm}
U desnoj karti susedna su i polja \(0\leftrightarrow2\) i \(4\leftrightarrow6\). Ovde zamišljamo presavijanje oko vertikalne ose i spajanje leve i desne ivice u prsten.
\end{frame}
